\documentclass{sigchi-ext}
% Please be sure that you have the dependencies (i.e., additional
% LaTeX packages) to compile this example.
\usepackage[T1]{fontenc}
\usepackage{textcomp}
\usepackage[scaled=.92]{helvet} % for proper fonts
\usepackage{graphicx} % for EPS use the graphics package instead
\usepackage{balance}  % for useful for balancing the last columns
\usepackage{booktabs} % for pretty table rules
\usepackage{ccicons}  % for Creative Commons citation icons
\usepackage{ragged2e} % for tighter hyphenation

% Some optional stuff you might like/need.
% \usepackage{marginnote} 
% \usepackage[shortlabels]{enumitem}
% \usepackage{paralist}
% \usepackage[utf8]{inputenc} % for a UTF8 editor only

%% EXAMPLE BEGIN -- HOW TO OVERRIDE THE DEFAULT COPYRIGHT STRIP --
% \copyrightinfo{Permission to make digital or hard copies of all or
% part of this work for personal or classroom use is granted without
% fee provided that copies are not made or distributed for profit or
% commercial advantage and that copies bear this notice and the full
% citation on the first page. Copyrights for components of this work
% owned by others than ACM must be honored. Abstracting with credit is
% permitted. To copy otherwise, or republish, to post on servers or to
% redistribute to lists, requires prior specific permission and/or a
% fee. Request permissions from permissions@acm.org.\\
% {\emph{CHI'14}}, April 26--May 1, 2014, Toronto, Canada. \\
% Copyright \copyright~2014 ACM ISBN/14/04...\$15.00. \\
% DOI string from ACM form confirmation}
%% EXAMPLE END

% Paper metadata (use plain text, for PDF inclusion and later
% re-using, if desired).  Use \emtpyauthor when submitting for review
% so you remain anonymous.
% \def\plaintitle{SIGCHI Extended Abstracts Sample File: Note Initial
%   Caps} \def\plainauthor{First Author, Second Author, Third Author,
%   Fourth Author, Fifth Author, Sixth Author}
\def\plaintitle{Do LLM Personas Decide Like Real Users? A Study of Value-Choice Coherence Across Expertise Levels}
\def\plainauthor{Rashi Tyagi, Jyoti Rawat, Ritik Hariani, Anil Muthigi, Saurav Nayak}
\def\plainkeywords{LLM personas; decision-making; value-choice alignment; expertise; human-AI comparison}
\def\plaingeneralterms{Human Factors, Experimentation}

\title{Do LLM Personas Decide Like Real Users? A Study of Value-Choice Coherence Across Expertise Levels}

\numberofauthors{5}
% Notice how author names are alternately typesetted to appear ordered
% in 2-column format; i.e., the first 4 autors on the first column and
% the other 4 auhors on the second column. Actually, it's up to you to
% strictly adhere to this author notation.
\author{%
  \alignauthor{%
    \textbf{Ritik Hariani}\\
    \affaddr{University of Illinois}\\
    \affaddr{Urbana-Champaign, IL, USA}\\
    \email{ritikh2@illinois.edu} }\alignauthor{%
    \textbf{Anil Muthigi}\\
    \affaddr{University of Illinois}\\
    \affaddr{Urbana-Champaign, IL, USA}\\
    \email{muthigi2@illinois.edu} } \vfil \alignauthor{%
    \textbf{Saurav Nayak}\\
    \affaddr{University of Illinois}\\
    \affaddr{Urbana-Champaign, IL, USA}\\
    \email{sgnayak2@illinois.edu} }\alignauthor{%
    \textbf{Jyoti Rawat}\\
    \affaddr{University of Illinois}\\
    \affaddr{Urbana-Champaign, IL, USA}\\
    \email{jrawat2@illinois.edu} } \vfil \alignauthor{%
    \textbf{Rashi Tyagi}\\
    \affaddr{University of Illinois}\\
    \affaddr{Urbana-Champaign, IL, USA}\\
    \email{rtyagi4@illinois.edu} } }

% Make sure hyperref comes last of your loaded packages, to give it a
% fighting chance of not being over-written, since its job is to
% redefine many LaTeX commands.
\definecolor{linkColor}{RGB}{6,125,233}
\hypersetup{%
  pdftitle={\plaintitle},
  pdfauthor={\plainauthor},
 % pdfauthor={\emptyauthor},
  pdfkeywords={\plainkeywords},
  bookmarksnumbered,
  pdfstartview={FitH},
  colorlinks,
  citecolor=black,
  filecolor=black,
  linkcolor=black,
  urlcolor=linkColor,
  breaklinks=true,
}


\begin{document}
%% For the camera ready, use the commands provided by the ACM in the Permission Release Form.
\CopyrightYear{2020}
\setcopyright{rightsretained}
\conferenceinfo{CHI'20,}{April  25--30, 2020, Honolulu, HI, USA}
\isbn{978-1-4503-6819-3/20/04}
\doi{https://doi.org/10.1145/3334480.XXXXXXX}
%% Then override the default copyright message with the \acmcopyright command.
\copyrightinfo{\acmcopyright}

\maketitle

% Uncomment to disable hyphenation (not recommended)
% https://twitter.com/anjirokhan/status/546046683331973120
\RaggedRight{} 

% Do not change the page size or page settings.
\begin{abstract}
Persona-conditioned LLMs are increasingly used as stand-ins for human participants in HCI research, yet little work has tested whether they replicate the decision-making patterns that actually distinguish real users. We ran a controlled comparison of value-choice coherence between 14 LLM personas (Gemini 2.5 Flash-Lite, 7 expert and 7 novice photographers, 10 runs each) and 14 screened human participants on a multi-attribute camera-purchasing task. All participants ranked 11 camera attributes by personal priority, then selected between 15 camera pairs presented at three complexity levels (3, 6, or 11 visible attributes). We measured alignment between stated priorities and actual choices, how that alignment changed with complexity, and how closely LLM and human patterns matched. On aggregate the two populations appear similar (mean similarity score = 0.953), but the cross-condition Pearson correlation is near zero (\textit{r} = 0.066), and the patterns diverge in direction. Human expert alignment degrades as complexity increases; LLM expert alignment improves. The expert-novice gap reverses sign at the highest complexity level. LLM novice personas are a reasonable proxy for human novices; LLM expert personas are not. These findings show that high-level outcome metrics can mask fundamentally different decision processes, and suggest that evaluating LLM-based synthetic users requires condition-level, process-oriented comparisons rather than aggregate alignment scores.
\end{abstract}



\vspace{0.8em}
\noindent\rule{0.42in}{0.4pt}\\[-0.2em]
{\footnotesize *All authors contributed equally to this work.}
\vspace{0.8em}


\keywords{\plainkeywords}

% ACM Classfication

\begin{CCSXML}
<ccs2012>
<concept>
<concept_id>10003120.10003121.10003122.10003334</concept_id>
<concept_desc>Human-centered computing~User studies</concept_desc>
<concept_significance>500</concept_significance>
</concept>
<concept>
<concept_id>10003120.10003121.10003122.10003333</concept_id>
<concept_desc>Human-centered computing~Empirical studies in HCI</concept_desc>
<concept_significance>300</concept_significance>
</concept>
<concept>
<concept_id>10010147.10010257.10010293</concept_id>
<concept_desc>Computing methodologies~Natural language processing</concept_desc>
<concept_significance>100</concept_significance>
</concept>
</ccs2012>
\end{CCSXML}

\ccsdesc[500]{Human-centered computing~User studies}
\ccsdesc[300]{Human-centered computing~Empirical studies in HCI}
\ccsdesc[100]{Computing methodologies~Natural language processing}


% Print the classficiation codes
\printccsdesc

\section{Introduction}
Large Language Models (LLMs) are being used more frequently in Human Computer Interaction (HCI) research to simulate human behavior, with applications extending over user feedback generation, usability evaluation, and early stage design testing \cite{hamalainen2023synthetichci, park2023generativeagents}. An increasing amount of work uses persona-conditioned LLMs as proxies for real users. These personas are assigned demographic attributes, domain expertise levels, or behavioral traits through prompting to approximate diverse populations at scale \cite{argyle2023out,aher2023usinglargelanguagemodels}. This approach has been prominent because it promises a cost-effective and rapid alternative to human participant recruitment process, which is particularly valued in interactive design contexts where fast feedback cycles are valuable. 

However, limited study has been done on the behavioral validity of LLM-based user simulation. Most evaluations assess whether synthetic outputs are plausible or useful to designers, rather than whether they reflect the underlying decision processes that characterize real human behavior \cite{hamalainen2023synthetichci, salminen2025syntheticusers}. Recent work has shown that LLMs can produce surface-level human-like responses while failing to capture deeper, context-dependent behavior \cite{santurkar}, and the persona grounding through prompting alone may not faithfully represent the cognitive and behavioral characteristics the persona claims to simulate \cite{aher2023usinglargelanguagemodels, shanahans}.Importantly, an LLM can give answers that look reasonable on the surface, even if the way it arrived at those answers is very different from how a human would. Most current evaluation methods do not capture that difference. 

This gap carries a direct consequence in HCI practice. Research in decision science has long established that human choices are shaped by task complexity, information availability, and domain expertise \cite{payne1993adaptive}. Experts tend to apply domain knowledge more consistently across decision contexts, whereas novices rely on surface-level cues or simplified heuristics as task complexity increases \cite{payne1993adaptive}. Value-choice inconsistency, where stated preferences, diverge from actual choices is a well-documented phenomenon in human decision making, especially under increasing cognitive load \cite{Tversky}. If persona-conditioned LLMs are used to generate synthetic evaluation data, they must not only produce plausible choices but replicate these characteristic patterns of human decision behavior. Whether they do so across varied complexity levels is an open empirical question. 

We address this question through a controlled comparison of value-choice coherence between persona conditioned LLMs and human participants in a multi-attribute camera purchasing task. We created 14 LLM personas using Gemini 2.5 Flash-Lite, of which 7 were experts and 7 were novices, each with detailed biographical backstories. Furthermore, we recruited 14 human participants screened to match the same expertise levels. All participants ranked 11 camera attributes by personal priority, then selected between camera pairs presented at three levels of complexity: simple (3 attributes), moderate (6 attributes), and complex (11 attributes). This design allows us to evaluate behavioral alignment across specific conditions where human decision patterns are known to shift. 

Our study is organized around two research questions, evaluated through three complementary metrics.\textbf{ RQ1 (Expertise and Alignment)}: Do expert LLM personas demonstrate higher value–choice alignment than novice personas, and does this difference reflect the alignment gap observed between human experts and novices?\textbf{ RQ2 (Complexity and Robustness)}: Does value-choice alignment decrease as product description complexity increases, and do expert personas exhibit greater robustness to this degradation compared to novice personas? To evaluate these questions, we measure three behavioral signals: \textbf{value–choice alignment} (whether choices match stated priorities), \textbf{consistency degradation} (whether alignment drops as decision complexity increases), and the \textbf{human–LLM persona gap} (whether LLM persona decision patterns match those of their human counterparts across expertise levels and complexity conditions).

Rather than evaluating synthetic users only on output plausibility alone, we evaluate them through decision-theoretic signals: how they reason across multiple attributes, how expertise changes their choices and whether their stated preferences stay consistent with their final decisions. Our findings reveal that while aggregate metrics suggest closeness between LLM and human alignment scores (similarity = 0.953), the cross-condition Pearson correlation is nearly zero, at r = 0.066, suggesting that the LLM scores do not track human scores in a consistent way. We also find that the direction of the complexity effect and the expert-novice gap differs substantially between the two populations. LLM expert personas improve in alignment as complexity rises, directly inverting the degradation pattern observed in the human experts. These results suggest that surface-level proxies such as choice accuracy or preference match should not alone be treated as evidence of behavioral validity, and motivates a process-level approach to evaluate LLM-based user simulation in HCI.

\section{Related Work}

\subsection{LLM-based User Simulation in HCI}

Simulation has long been used in HCI to explore user behavior and support early-stage design decisions \cite{murraysmith2022simulation}. With the rise of large language models (LLMs), recent work has extended this paradigm toward generating synthetic user data. H\"am\"al\"ainen et al.\ show that LLMs can produce realistic user-generated content for ideation and piloting, but raise concerns about ecological validity \cite{hamalainen2023synthetichci}. Similarly, Xiang et al.\ introduce SimUser, which simulates diverse users to generate usability feedback at scale \cite{xiang2024simuser}. 

A parallel line of work explores LLM-based agents as models of human-like behavior. Park et al.\ demonstrate that generative agents can produce coherent, interactive behaviors over time \cite{park2023generativeagents}. However, these approaches primarily evaluate synthetic users based on plausibility and usefulness of generated outputs rather than whether they reflect underlying human decision processes. Recent work further shows that LLM-simulated users can be unreliable proxies for real users, particularly in behavioral evaluations rather than surface-level tasks \cite{lostinsimulation2024}. Salminen et al.\ further highlight that LLM-based synthetic users may be limited by bias and lack of grounding in real human experience \cite{salminen2025syntheticusers}.

\subsection{Decision Behavior, Expertise, and Preference Consistency}

Research in decision-making shows that human choices are shaped by task complexity, information availability, and expertise. In multi-attribute decision contexts, individuals adapt their strategies as complexity increases, often shifting from systematic trade-offs to simpler heuristics \cite{payne1993adaptive}. In such settings, individuals evaluate tradeoffs between competing attributes, often prioritizing some features over others depending on task complexity. Expertise plays a critical role: experts tend to apply domain knowledge more consistently, while novices rely on surface-level cues or simplified strategies \cite{chi1988nature, shanteau1992competence}. 

Prior work also shows that alignment between stated preferences and actual choices is not always stable. Preference reversals demonstrate that individuals may express one preference but act differently when making decisions, particularly under varying decision frames or information conditions \cite{slovic1983preference}. These findings suggest that value--decision coherence is sensitive to both cognitive load and expertise.

Recent work in human-AI interaction further shows that decision outcomes are influenced by how users interpret and integrate model outputs \cite{bansal2021does}. At the same time, LLMs can approximate certain human response patterns but often fail to capture deeper, context-dependent decision behavior \cite{argyle2023out}.

\subsection{Our Contribution}

Our work bridges these areas by evaluating whether persona-conditioned LLMs exhibit human-like decision behavior rather than merely producing plausible outputs. We draw on decision-theoretic concepts such as multi-attribute reasoning, expertise effects, and preference-choice consistency to design controlled comparisons between human participants and LLM personas.

Specifically, we compare decision patterns across expert and novice conditions while varying the amount of decision-relevant information. This allows us to test whether LLM personas exhibit the same breakdowns and adaptations in value--decision coherence observed in human decision-making. By grounding synthetic user evaluation in comparative decision behavior, our work provides a more rigorous framework for assessing the validity of LLM-based user simulation in HCI.

\section{Methodology}

\subsection{Overview}
This section describes the participant groups, the camera-purchasing task used as the experimental domain, how sessions were structured, the three metrics used to evaluate behavior, and the approach taken to analyze the results.

\subsection{Participants}

\subsubsection{LLM Personas}
We created 14 personas and ran them on Gemini 2.5 Flash-Lite, 7 experts and 7 novices. The expert group covered a range of professional photography contexts: wedding and portrait work, wildlife, street and documentary, commercial product, landscape and astrophotography, fashion editorial, and sports. The novice group was equally varied: a college student, a dad who photographs family events, a food blogger, a retired teacher just getting into photography, a travel enthusiast, a college athlete making fitness content, and an event planner.

Rather than just labeling a persona ``expert'' or ``novice'' and telling it what to care about, each persona had a detailed backstory covering gear they own, how long they have been shooting, what they know and don't know, and what they are trying to buy and why. The idea was that their priorities should come through naturally from the context, not because they were told to prioritize certain things. The prompts were fixed and never changed between runs. We ran each persona 10 times to capture behavioral variance across generations rather than rely on a single response, giving us 140 sessions in total. All sessions were run using the Gemini 2.5 Flash-Lite API with temperature 0.75, top-p 0.95, and a maximum of 1024 output tokens per turn. Temperature was kept above zero intentionally so that responses varied meaningfully across runs.

\subsubsection{Human Participants}
For the human side, we recruited 14 people, 7 experts and 7 novices, through photography clubs and personal contacts. Each person was approached individually, informed about the study, and asked to participate voluntarily. There was no compensation. Screening was done in person by asking each candidate the qualifying questions directly: experts had to own at least three cameras, have bought gear in the last two years, and be comfortable explaining technical concepts like aperture and ISO; novices had to be phone-only shooters who had never bought a dedicated camera. The same screening questions were also included at the start of the Google Form as a self-report and for documentation purposes. Once screened in, participants were given written instructions walking them through the task, then completed the form anonymously at their own pace. There were no practice rounds.

\subsection{Stimuli}
We worked with a pool of 11 camera attributes: price, image quality, battery life, optical zoom, weight, ease of use, manual controls, RAW support, burst speed, weather sealing, and 4K video. All camera names and specs were generated synthetically using Claude. Names were made up (e.g., ``Alvex T1,'' ``Brikon S2'') so no one could rely on brand familiarity, and attribute values were set carefully so that each pair had no obvious winner: every choice required a genuine tradeoff between attributes the participant actually cared about \cite{LLM-Based_User_Simulation_2026}.

Complexity was defined by how many attributes were visible for a given pair:

\begin{table}[h]
  \centering
  \begin{tabular}{l r}
    \toprule
    \textbf{Complexity} & \textbf{Visible Attributes} \\
    \midrule
    Simple   & 3 \\
    Moderate & 6 \\
    Complex  & 11 (all) \\
    \bottomrule
  \end{tabular}
  \caption{Number of visible attributes per complexity level.}
\end{table}

There were 5 pairs at each complexity level, so 15 pairs per session total. Pairs within the same complexity level did not all show the same attributes. One simple pair might show price, image quality, and battery life, while another showed weight, ease of use, and manual controls. This ensured no single tradeoff dominated the design and that participants had to think carefully about different combinations of features at every level.

\subsection{Session Structure}
Each session covered all 15 pairs in one go. For LLM personas, this was a single multi-turn conversation. For human participants, the same structure was followed through the Google Form.

\textbf{Step 1: Persona setup.}
For LLM sessions, the persona's biographical prompt was loaded at the start and kept the same throughout. Nothing was edited mid-session.

\textbf{Step 2: Priority ranking.}
Before seeing any cameras, participants ranked all 11 attributes from most to least important. This ranking came from them directly before any cameras were shown, and it was the same ranking used to evaluate all 15 pairs later.

\textbf{Step 3: Camera pairs.}
Pairs were shown one at a time in a fixed order: 5 simple pairs first, then 5 moderate, then 5 complex. Only the attributes belonging to that pair's subset were shown, written out in plain language rather than as raw variable names. After seeing each pair, participants picked one camera by submitting: \texttt{Choice: [camera name]}.

The order never changed across sessions, so any effects from seeing earlier pairs carry through consistently for everyone. Human participants followed the same structure through the Google Form, with the priority ranking presented first and the 15 camera pairs following in the same fixed order.

\subsection{Evaluation Metrics}

\subsubsection{Value Choice Alignment}
For each pair, we first filtered down to the visible attributes where the two cameras actually differed. Those differing attributes were then sorted by the participant's priority ranking, most important first, and assigned local ranks 1 through $n$. A linear weight was applied to each: the top-ranked differing attribute received weight $n$, the next received $n-1$, and so on down to 1 for the least important. For each attribute, we checked whether the chosen camera was objectively better (based on known attribute directionality, for example lower price is better, higher image quality is better), scoring it 1 if correct and 0 otherwise. The pair score was the weighted sum of these correctness values divided by the total weight, giving a number between 0 and 1.

These pair scores were then aggregated upward. The five pair scores within a session and complexity level were summed to give a session score between 0 and 5. For LLM personas, session scores were averaged across the 10 runs per persona to give a persona score. Finally, persona scores were averaged across the 7 personas in each group and divided by 5 to produce a final alignment score between 0 and 1 for each label and complexity combination.

\subsubsection{Consistency Degradation under Feature Complexity}
Rather than fitting a regression across all three levels, we measured alignment change at two specific transitions: simple to moderate and moderate to complex. For each session, the degradation at each step was the difference in alignment score between adjacent levels. A negative value meant alignment dropped at that transition; a positive value meant it improved. These session-level differences were averaged across the 10 runs per LLM persona, then averaged again across the 7 personas per group, giving a mean transition slope for experts and novices separately.

\subsubsection{Human-LLM Persona Gap}
For each of the six label and complexity combinations, we computed the signed gap between LLM and human alignment scores (LLM minus human) and the absolute gap. These were summarized in two ways: per label averaged across complexity levels, and the expert minus novice gap within each population across complexity levels. We also computed a similarity score defined as 1 minus the mean absolute gap across all six conditions, and a Pearson correlation between LLM and human scores across those conditions. The correlation captures whether LLM and human scores follow the same pattern across conditions, not just whether they land in the same numerical range.

\subsection{Statistical Analysis}
All analysis was descriptive and computed directly from the pairwise alignment scores. For value choice alignment, we computed means and standard deviations at the group and persona level across all three complexity conditions. For consistency degradation, we reported the two pairwise transition differences (simple to moderate and moderate to complex) per group. For the human-LLM gap, we reported mean signed and absolute gaps per condition, the similarity score, and the Pearson correlation between LLM and human scores across all six label and complexity combinations. No inferential model was fit; conclusions were drawn from the direction and magnitude of these descriptive comparisons.

\section{Results}
We report results across three metrics: value-choice alignment (VCA), consistency degradation across complexity levels, and the human-LLM persona gap.

\subsection{Value-Choice Alignment}
Table~\ref{tab:vca} summarises mean VCA scores for each (label $\times$ complexity) condition. LLM experts score 0.506 at simple, 0.558 at moderate, and 0.605 at complex. LLM novices score 0.657 at simple, 0.541 at moderate, and 0.558 at complex. Human experts score 0.590, 0.527, and 0.518 across the same levels; human novices score 0.590, 0.535, and 0.549. Across complexity levels, LLM personas average 0.556 (expert) and 0.585 (novice); humans average 0.545 (expert) and 0.558 (novice).

\begin{table}[h]
  \centering
  \begin{tabular}{l l r r}
    \toprule
    \textbf{Label} & \textbf{Complexity} & \textbf{LLM} & \textbf{Human} \\
    \midrule
    Expert & Simple   & 0.506 & 0.590 \\
    Expert & Moderate & 0.558 & 0.527 \\
    Expert & Complex  & 0.605 & 0.518 \\
    Novice & Simple   & 0.657 & 0.590 \\
    Novice & Moderate & 0.541 & 0.535 \\
    Novice & Complex  & 0.558 & 0.549 \\
    \bottomrule
  \end{tabular}
  \caption{Mean value-choice alignment by population, label, and complexity.}~\label{tab:vca}
\end{table}

\subsection{Consistency Degradation}
For human experts, alignment drops at both transitions (simple $\rightarrow$ moderate $= -0.064$, moderate $\rightarrow$ complex $= -0.008$), consistent with degradation as more attributes are introduced. Human novices show a similar drop simple $\rightarrow$ moderate ($-0.056$) but a small recovery moderate $\rightarrow$ complex ($+0.015$). LLM experts show the opposite pattern: alignment \emph{improves} at both transitions (simple $\rightarrow$ moderate $= +0.052$, moderate $\rightarrow$ complex $= +0.047$). LLM novices degrade sharply simple $\rightarrow$ moderate ($-0.116$) before stabilising moderate $\rightarrow$ complex ($+0.018$).

\subsection{Human-LLM Persona Gap}
Across all six (label $\times$ complexity) conditions, the mean absolute gap between LLM and human VCA scores is 0.047 and the mean signed gap is $+0.019$, yielding a similarity score of 0.953. However, the Pearson correlation between LLM and human scores across the six conditions is 0.066, effectively zero. Per-condition signed gaps (LLM minus human) are largest at expert/simple ($-0.085$) and expert/complex ($+0.086$), and smallest at novice/moderate ($+0.006$) and novice/complex ($+0.009$). The expert-novice gap within each population also diverges by complexity: at simple, LLM novices outscore LLM experts by 0.151 while humans are tied (0.0); at complex, LLM experts outscore LLM novices by 0.046 while human experts score 0.031 \emph{below} human novices, a reversal in direction.

\subsection{Key Findings}

\textbf{Surface similarity is misleading.}
The high similarity score (0.953) and small mean absolute gap (0.047) suggest LLM and human alignment scores are numerically close. The near-zero Pearson correlation (0.066) shows they do not track the same pattern across conditions; the populations land in the same numerical range without responding to expertise or complexity in the same way.

\textbf{LLM experts invert the human complexity trend.}
Human experts score highest at simple (0.590) and lowest at complex (0.518), consistent degradation with complexity. LLM experts do the opposite: lowest at simple (0.506), highest at complex (0.605). This is the sharpest divergence in the dataset and runs counter to the expected expertise effect.

\textbf{LLM novices track humans more closely than LLM experts do.}
Novice gaps at moderate ($+0.006$) and complex ($+0.009$) are negligible, and both groups share the simple $\rightarrow$ moderate drop. The dominant LLM-novice deviation is the unusually high simple-condition score (0.657 vs.\ 0.590). Overall, LLM novice personas are a reasonable proxy for human novices; LLM expert personas are not.

\textbf{The expert-novice gap reverses direction at complex.}
Within humans, novices slightly outperform experts at complex (gap $= -0.031$). Within LLMs, experts outperform novices at the same level (gap $= +0.046$). The direction is reversed; the LLM expert persona does not replicate expert human behavior.

\textbf{Largest single mismatch: expert/simple (gap=$-0.085$).}
LLM experts score 8.5\% points below human experts at simple complexity. At the same level, human experts and novices are tied (0.590 each), while LLM novices massively outscore LLM experts (0.657 vs.\ 0.506).

\textbf{Aggregate metrics can mask divergent decision processes.}
A single-number similarity score would have characterized the LLM and human populations as well aligned. The condition-level breakdown and the cross-condition correlation reveal that this agreement is largely coincidental, motivating process-level rather than outcome-level evaluation of LLM-based user simulation (see Future Work).



\section{Conclusion}
Our two research questions are evaluated through three behavioral metrics: value-choice alignment, consistency degradation, and the human-LLM persona gap, and the pattern across all three converges on the same conclusion.

On \textbf{RQ1 (Expertise and Alignment)}, expert LLM personas do not demonstrate higher value-choice alignment than novice personas, they actually scored lower on average than novice personas (0.556 vs. 0.585). More critically, this gap does not reflect the human pattern. Human experts and novices are tied at simple complexity (0.590 each), and human experts are slightly behind novices at complex (0.518 vs. 0.549, gap = -0.031). LLM personas show a different picture: at simple complexity, LLM novices massively outscore LLM experts (0.657 vs. 0.506), and at complex, LLM experts lead novices (0.605 vs. 0.558, gap = +0.046), a path inverse to that of humans. Therefore, Labeling a persona as an "expert" through prompting does not result in expert-like decision making. Notably, this failure is asymmetric: LLM novice personas are a reasonable proxy for human novices, particularly at moderate and complex conditions where gaps are negligible (+0.006 and +0.009 respectively); LLM expert personas are not.

On \textbf{RQ2 (Complexity and Robustness)}, the populations diverge sharply and in opposite directions. Human experts show the expected degradation as complexity rises (simple -> moderate = -0.064, moderate -> complex = -0.008). LLM experts show the opposite: alignment improves at both transitions (+0.052, +0.047), directly inverting the human degradation pattern. LLM experts do not exhibit greater robustness to complexity in any human-like sense. They exhibit a fundamentally different response to it. LLM novices, by contrast, degrade sharply from simple to moderate (-0.116) before stabilizing (moderate -> complex = +0.018), which at least shares the direction of human novice degradation (-0.056), even if the magnitude differs.

The human-LLM persona gap metric shows why this matters. Aggregate scores look fine, with a similarity of 0.953 and a mean absolute gap of just 0.047. But the cross-condition Pearson correlation is near zero (r = 0.066) which means the two populations land in the same numerical range by coincidence, not because they are responding to expertise and complexity the same way. The largest gaps fall exactly where expert behavior diverges most, at simple (-0.085) and complex (+0.086) complexity.

These findings suggest that current persona prompting does not ground ``expertise'' in the way that produces human like decision behavior, and that aggregate alignment metrics can obscure this. For HCI work using LLM-based synthetic users, the implication is that outcome-level proxies (e.g., choice accuracy, preference match) should not by themselves be treated as evidence of behavioral validity; researchers should evaluate whether synthetic users replicate the \emph{patterns} of degradation, expertise, and trade-off that motivate the simulation in the first place. Process-level validation, scaled across model families and decision domains, is the more demanding bar this work points toward.

\section{Limitations}
Several limitations should be considered when interpreting these findings.
First, our human sample is small ($N = 14$; 7 experts, 7 novices), recruited through convenience sampling without compensation. Statistical power is limited and the participant pool may be biased toward photography enthusiasts.

Second, all LLM data come from a single model (Gemini 2.5 Flash-Lite) under a single prompting configuration. Behavior may differ substantially across model families, sizes, or providers; our results should not be read as a general claim about ``LLM personas,'' only about this model in this configuration.

Third, the camera stimuli were generated by a different LLM (Claude) than the one being tested. This avoids hand-curated stimuli leaking experimenter  expectations, but introduces a subtle distributional alignment between stimuli and at least one model family that could affect cross-model replication.
      
Fourth, we use a single, low-stakes decision domain (consumer camera purchasing). Whether the inverted complexity trend in LLM experts generalizes to higher-stakes domains where novice/expert differences are more consequential (medical, financial, legal) is an open question this work cannot answer.

Fifth, we report descriptive statistics only and do not fit a mixed-effects or other inferential model. The directional effects we highlight, particularly the inverted complexity trend in LLM experts and the sign reversal of the expert-novice gap at complex, should therefore be treated as suggestive rather than confirmatory.

Sixth, the priority ranking is treated as a stable preference, but participants completed it once without practice rounds and may have ranked differently with reflection. The alignment scores are sensitive to this single-shot ranking.      

Finally, pair order was fixed (simple $\rightarrow$ moderate $\rightarrow$ complex) for both populations. This controls between-condition order effects but bundles session-internal learning, fatigue, and anchoring into the complexity manipulation, so they cannot be separated from it.

\section{Future Work}
These findings motivate several directions.

First, the divergent complexity trends in RQ2 are descriptive. The natural confirmatory test is a mixed-effects logistic regression with persona label, complexity, and their interaction as fixed effects, and persona nested within label as a random effect. This requires enlarging the human sample; until that is done, the inverted trend in LLM experts remains a strong suggestion rather than an established effect. 

Second, our results are anchored to a single model. Replicating the same protocol across model families (GPT-class, Claude-class, open-weight models) and sizes will determine whether value--choice degradation under complexity, and its inversion in LLM experts, is a property of one model or a general property of persona-conditioned LLMs. The human side requires no changes.

Third, the contrast between high aggregate similarity ($0.953$) and near-zero per-condition correlation ($r = 0.066$) is itself a methodological finding. HCI evaluation of LLM-based synthetic users needs metrics that capture decision \emph{processes} (which attributes drive each choice, how priority is traded against attribute values, where the model deviates from its own stated ranking), rather than only outcome-level alignment. Developing and validating such metrics is, in our view, the most important next step for the validity of LLM persona research.

Fourth, our domain has low decision stakes. Extending the protocol to higher-stakes domains, where expert/novice differences carry real consequences, will test whether the LLM-expert-improving-with-complexity pattern is general or an artifact of low-stakes consumer choice. This is also the setting where persona validity matters most. 

A fifth, more speculative direction addresses RQ1 directly: whether richer expertise grounding (retrieval over domain knowledge, few-shot expert traces, or chain-of-thought elicitation of expert reasoning) can close the gap between LLM and human expert behavior

\balance{} 

\bibliographystyle{SIGCHI-Reference-Format}
\bibliography{sample}

\end{document}
